\section{哈密顿作用量原理}
\label{sec:hamilton-action-principle}

牛顿力学中的拉格朗日量具有形式
\begin{equation}\label{eq:newtonian-lagrangian}
  L=T-U,
\end{equation}
其中 $T$ 是\emph{动能}
\[
  T=T(x,\dot x)=\frac12m\dot x\cdot\dot x,
\]
而 $U=U(x):\mathbb R^n\to\mathbb R$ 是\emph{势能}，它是只依赖位置向量 $x\in\mathbb R^n$ 的函数。于是牛顿运动方程
\begin{equation}\label{eq:newton-motion}
  \frac d{dt}(m\dot q^i)=-\frac{\partial U}{\partial q^i},
  \qquad i=1,\ldots,n,
\end{equation}
等价于拉格朗日量~\eqref{eq:newtonian-lagrangian} 的欧拉--拉格朗日方程~\eqref{eq:euler-lagrange}。由此得到：

\medskip
\noindent\textbf{哈密顿最小作用量原理。}
\stfhSubIndex{I051}{zuo^yong^liang^fan^han}{作用量泛函（action functional）}{ha^mi^dun^zuo^yong^liang^fan^han}{哈密顿作用量泛函（Hamilton’s action functional）}
服从牛顿第二定律的力学系统运动，恰好是泛函
\[
  \Phi(\gamma)=\int_{t_0}^{t_1}L\,dt
\]
的极值道路（\emph{边界条件应取前述适当形式}）。

回顾牛顿力学中\emph{动量向量}的定义
\begin{equation}\label{eq:newtonian-momentum}
  m\dot x\coloneqq p,
\end{equation}
而\emph{力场}由
\[
  -\nabla U\coloneqq F
\]
给出。由于式~\eqref{eq:newtonian-lagrangian} 中拉格朗日量的特殊形式，$p_i=m\dot q^i$ 恰好就是 $\partial L/\partial\dot q^i$。这引出经典力学中的下述概念。（更多讨论见 (Go80) 或 (Ar89)。）

\begin{definition}
\label{def:generalized-momenta}
设 $L$ 是任意拉格朗日量。\emph{广义动量}记为 $p_i$，定义为
\begin{equation}\label{eq:generalized-momenta}
  p_i=\frac{\partial L}{\partial\dot q^i}.
\end{equation}
\end{definition}

在此定义下，欧拉--拉格朗日方程变为
\begin{equation}\label{eq:momentum-euler-lagrange}
  \dot p_i=\frac{\partial L}{\partial q^i},
  \qquad i=1,\ldots,n.
\end{equation}

\begin{exercise}
\label{ex:generalized-momenta-invariant}
以不依赖坐标的方式解释式~\eqref{eq:generalized-momenta} 和~\eqref{eq:momentum-euler-lagrange}。说明为什么应把它们的右端看作协变一阶张量，或微分 $1$-形式。
\end{exercise}
